\section{Method: \model}
Employing the data from \mind, we introduce an exploratory framework, \model, for our task, leveraging the power of LLMs.
Raw HTML documents, which could consist of thousands of elements, are either infeasible or cost-prohibitive to be directly fed into LLMs.
We propose a two-stage process that synergizes the strength of small and large LMs, as shown in Figure~\ref{fig:overview}.
In the first stage, a fine-tuned small LM is used to rank the elements present on a webpage, yielding a small pool of promising candidates.
In the second stage, these candidate elements are consolidated to form a representative snippet of the webpage, which is then processed by an LLM to predict the final action, including predicting both the element for interaction and the corresponding operation.

\subsection{Candidate Generation with Small LMs}
\begin{figure}
    \begin{minipage}{.52\linewidth}
       \centering
        \includegraphics[width=0.9\textwidth]{figures/pipeline.pdf}
        \caption{The overall pipeline for \model\ with a small ranking LM for candidate generation, and a large prediction LM for action prediction.}
        \label{fig:overview}
    \end{minipage}%
    \hfill
    \begin{minipage}{.44\linewidth}
        \centering
        \includegraphics[width=0.85\textwidth]{figures/model_0.pdf}
        \caption{Illustration of the candidate generation module and the templates for constructing task query and candidate representation.}
        \label{fig:candidate_generation}
    \end{minipage} 
    \vspace{-1em}
\end{figure}
Given the task description, the snapshot of the webpage at step $t$, and the actions performed in the preceding $t-1$ steps, we treat candidate generation as a ranking task. The task is to select the top-$k$ candidate DOM elements from the webpage that best align with both the task description and the current step.
We formulate the task query by concatenating the task description with previous actions.
The textual representation of each candidate DOM element is derived from a combination of the element's tag, its textual content, and salient attribute values, as well as the textual representation of its parent and child elements.
As shown in Figure~\ref{fig:candidate_generation}, we pair each DOM element with the task query and feed it to an encoder-only LM through the cross-encoder architecture~\cite{reimers-2019-sentence-bert}, yielding a matching score.
At training time, we randomly sample negative elements from the webpage, and use the target element as the positive example. The matching score is passed through a sigmoid activation function and optimized with a binary cross entropy loss. At inference time, we score all elements in the webpage and pick the top-$k$ elements with the largest logits as input to the second stage.

\subsection{Action Prediction with LLMs}
After obtaining the top-$k$ candidates, we utilize the candidate set to prune the webpage snapshot and construct snippets that only include the selected candidates and their neighbours as inputs to an LLM.
Recent studies~\cite{flant5,pangu} have suggested that training LMs for discrimination rather than generation is more generalizable and sample-efficient for other grounding tasks.
Inspired by that, we convert the task of element selection into a \textit{multi-choice question answering} (QA) problem. 
Instead of generating the complete target element, the LM is trained to instead select from a list of options.
For comparison, we also include a baseline that directly generates the target element based on the provided webpage snippet.
In both cases, we directly let the LLM generate the operation, along with the additional value needed for some operations.
An example is shown in Figure~\ref{fig:template}.
We incorporate up to \num{5} candidate elements within each input, together with a None option, and partition the candidate set into several groups.
During training, we construct the target sequence using ground-truth actions and fine-tune the model using a left-to-right language modeling objective.
During inference, we divide the top-$k$ candidates into multiple clusters of five options. If more than one option is selected after a round, we form new groups with the selected ones. This process repeats until a single element is selected, or all options are rejected by the model, i.e., the model chooses the None option for all groups.
\begin{figure}
    \centering
    \includegraphics[width=0.9\textwidth]{figures/model_1.pdf}
    \caption{Illustration of action prediction with LLMs.}
    \label{fig:template}
    \vspace{-1.5em}
\end{figure}
